\documentclass[10pt,a4paper]{amsart}
\usepackage[margin=19mm]{geometry}
\usepackage{amsmath,amssymb,mathtools}
\usepackage[scr=rsfs,cal=euler]{mathalfa}
\usepackage{xfrac}
\usepackage[unicode,hidelinks,bookmarks=false]{hyperref}
\newcommand{\ie}{i.e.\ }
\newcommand{\cf}{cf.\ }
\newcommand{\resp}{respectively\ }
\newcommand{\Subsection}{\subsection}
\providecommand{\nobreakdash}{\nobreak\hbox{-}}
\setlength{\emergencystretch}{2em}
\multlinegap0pt
\usepackage{amssymb}
\usepackage[normalem]{ulem}
\usepackage{tikz}
\newtheorem{theorem}{Theorem}
\title{Structural and Spectral Properties of Strictly Interval Graphs}
\author{Claudia Justel, Lilian Markenzon}
\date{Source: arXiv:2512.01175v2 (CC BY 4.0)}
\begin{document}
\maketitle

\noindent\emph{Source and licence.} Structural and Spectral Properties of Strictly Interval Graphs. Version arXiv:2512.01175v2 (CC BY 4.0), distributed by arXiv under the Creative Commons Attribution 4.0 International licence, http://creativecommons.org/licenses/by/4.0/, as recorded in the arXiv licence metadata for this exact version and shown on the abstract page https://arxiv.org/abs/2512.01175v2. Full licence notice: \texttt{licenses/arxiv-2512.01175-CC-BY-4.0.txt}.\par\smallskip

\noindent\textit{Editorial scope.} Counted unit: the article's theorem theo:counting (source0.tex lines 630--633). The article's complete original proof of it is delivered with it (source0.tex lines 634--642, one \texttt{proof} environment of 9 lines). No mathematics is written, repaired or re-derived: the statement and the proof are the article's own lines, in the article's own notation, and no retained line is substituted or re-flowed.  No further source-local context is needed: every cross-reference of every retained line resolves inside the two delivered spans.  Nothing was omitted from any retained span, so the package carries no omission record.  No retained line contains a citation, so the wrapper carries no bibliography and no bibliography entry.  No retained line contains a figure environment, an image-inclusion command or drawing code, so no image is delivered and the package declares no asset dependency.  Editorial formatting only, in addition: an AMS \texttt{amsart} preamble that restates the article's own declarations and macros used by the retained lines, namely the environments \texttt{theorem} on separate counters, together with the standard packages those lines need.  The retained environments print in this document as Theorem 1 in order of appearance, whereas the article prints the same items with the numbering of its own full document.  The complete native source of the article as distributed in the arXiv e-print for this exact version is bundled unchanged as \texttt{sources/190107/source0.tex}.
\begin{theorem}\label{theo:counting}
Let $s,p$ be fixed integer values.
There are $\frac{a(p)(a(p)+1)}{2}$ non-isormophic $SI$-core graphs.  
\end{theorem}

\begin{proof}
Let $G \in {\mathcal G}(s,p)$.
Each partition of the $p$ simplicial vertices adjacent to a $mvs$ can be combined
with a partition of the other $p$ simplicial vertices.
Hence, with the first partition we have $p$ diferent combinations.
With the second partition we have $p-1$ different combinations  and so on.
So, the number of combinations without repetition is given by 
the sum of the elements of an arithmetic progression from 1 to  $p$. 
\end{proof}

\end{document}
